\documentclass[10pt,a4paper]{amsart}
\usepackage[margin=19mm]{geometry}
\usepackage{amsmath,amssymb,mathtools}
\usepackage[scr=rsfs,cal=euler]{mathalfa}
\usepackage{xfrac}
\usepackage[unicode,hidelinks,bookmarks=false]{hyperref}
\newcommand{\ie}{i.e.\ }
\newcommand{\cf}{cf.\ }
\newcommand{\resp}{respectively\ }
\newcommand{\Subsection}{\subsection}
\providecommand{\nobreakdash}{\nobreak\hbox{-}}
\setlength{\emergencystretch}{2em}
\multlinegap0pt
\usepackage{mathtools}
\usepackage{physics}
\usepackage{enumitem}
\usepackage{multicol}
\usepackage{tcolorbox}
\usepackage{dsfont}
\usepackage{tikz-cd}
\usetikzlibrary{backgrounds}
\usepackage{longtable}
\usepackage{bbm}
\def\blfootnote{\gdef\@thefnmark{}\@footnotetext}
\newtheorem{theorem}{Theorem}
\newtheorem*{theorem*}{Theorem}
\newtheorem{corollary}{Corollary}[theorem]
\newtheorem{conj}[theorem]{Conjecture}
\newtheorem{lemma}[theorem]{Lemma}
\newtheorem{defn}[theorem]{Definition}
\newtheorem{propn}[theorem]{Proposition}
\newenvironment{sproof}{%
  \renewcommand{\proofname}{Sketch Proof}\proof}{\endproof}
\newcommand{\B}[1]{\mathbb{#1}}
\newcommand{\C}[1]{\mathcal{#1}}
\newcommand{\oop}{\preccurlyeq}
\newcommand{\is}{\mathbbm{1}}
\renewcommand{\arraystretch}{1.5}
\usepackage{parskip}
\newcommand{\JB}[1]{\textcolor{blue!70!black}{\textbf{JB:} #1}}
\begin{document}
\title[Branching-selection particle systems and inverse first passa]{Branching-selection particle systems and inverse first passage problems}
\author{Jacob Mercer}
\date{}
\maketitle
\noindent\emph{Source and licence.} Branching-selection particle systems and inverse first passage problems. Version arXiv:2606.13487v1 (CC BY 4.0). This material is distributed under the Creative Commons Attribution 4.0 International licence, http://creativecommons.org/licenses/by/4.0/, as recorded in the arXiv OAI licence metadata for this exact version and shown on the abstract page https://arxiv.org/abs/2606.13487v1. Full rights notice: licenses/arxiv-2606.13487-CC-BY-4.0.txt.\par\smallskip
\noindent\textit{Editorial scope.} Counted unit: the original \texttt{theorem} environment \texttt{2ndMainThemOmega} at source lines 149--155 together with its complete proof at lines 157--178; the delivered document keeps source lines 99--179 so that every referenced label and every citation resolves inside it. Native editable source is the paper's own concatenated TeX; the AMS wrapper is editorial formatting only. No publisher class, upstream script or unrelated artwork is redistributed.
\par\smallskip\noindent\textit{Reference note.} This article ships no compiled bibliography (biblatex); the entries delivered here are composed from the author's own BibTeX (.bib) fields, with no field invented and no entry dropped.
\begin{align} \label{omegaHydroLim}
    \begin{cases}
    u_t(x,t)=\frac{1}{2}u_{xx}(x,t) + f(t)u(x,t) - \Lambda\omega(x-b(t))u(x,t) & x\in \B{R}, t>0, \\
    u(x,0)=\rho(x) & x\in \B{R}\\
    \int_\B{R} u(y,t)dy=1\\
    %\int_{\B{R}}\omega(y-b(0))u(y,0)dy = f(0)
\end{cases}
\end{align}
and show that if $(u,b)$ is a solution, then $b$ solves the SKIFPP with initial distribution $\rho$ and $p([t,\infty))=\exp(-\int_0^t f(s)ds)$. In the particle system, particles branch at rate $f$ and a particle at location $x$ is deleted with probability proportional to $\omega(x-b^N(t))$ where $b^N(t)$ is a function of the empirical distribution at time $t$. Accordingly, $b^N(t)$ can be seen as an approximation of the solution $b$ of the SKIFPP. 

In the specific case that $\omega(x)=\is_{x\leq 0}$, if we choose $f(t)\equiv 1$ and $\Lambda\in (1,\infty)$, then we can combine the PDE and the integral condition to yield the non-linear PDE
\begin{align*}
    u_t(x,t)=\frac12 u_{xx}(x,t)+u(x,t)\left(1-\Lambda\is_{\int_{-\infty}^x u(y,t)dy \leq \Lambda^{-1}}\right).
\end{align*}
In this case, $b(t)$ is the point left of which a proportion $\Lambda^{-1}$ of the mass of $u$ lies. %In terms of the particle system, $(\is_{x\leq 0},1,\alpha^{-1},N)$-BBM is the branching Brownian motion in which at each branching event, we choose a particle uniformly from among the $\lfloor \alpha N\rfloor$ leftmost to be deleted. 

\section{Construction of the process}\label{constructionSec}

Essentially this process will be a branching Brownian motion in which a particle at location $X(t)$ is deleted at rate $\Lambda\omega(X(t)-b^N(t))$ where $b^N:[0,\infty)\to \B{R}$ is a reference point chosen so that the total branching rate at time $t$ is $Nf(t)$. Branching and deletion will happen simultaneously so that there are always exaclty $N$ particles. 

We now give a formal construction. Throughout this work, we will assume that $\omega:\B{R}\to [0,1]$, the \textit{selection weighting}, is a monotone decreasing function with $\omega(x)\to 1$ as $x\to-\infty$ and $\omega(x)\to 0$ as $x\to\infty$. We also assume that $f:[0,\infty)\to [0,\infty)$, the \textit{branching rate} function,is bounded and that $\Lambda > \|f\|_\infty$. 

Then we will describe the $(\omega,f,\Lambda,N)$-BBM at time $t$ by the probability measure $\mu_t^N:=\frac1N\sum_{i=1}^N \delta_{X_i^N(t)}$ where $X^N(t):=(X_1^N(t),\ldots,X_N^N(t))$ denote the locations of the $N$ particles in the system. Let $X^N(0)\in \B{R}^N$ be some initial configuration with $X_1^N(0)\leq X_2^N(0)\leq \cdots\leq X^N_N(0)$ and define $\rho^N:=\frac1N\sum_{i=1}^N \delta_{X^N_i(0)}$. Suppose that $\rho^N$ converges weakly to $\rho$, with $\rho\ll\lambda$ ($\rho$ is absolutely continuous with respect to the Lebesgue measure $\lambda$). We will use $\rho$ to denote both the probability distribution and its density function.

Let $\C{N}(t)$ be a Poisson process of intensity $Nf(t)$ and let the discontinuities of $\C{N}(t)$ be $\tau_1<\tau_2<\cdots$. These will be the branching times of the process. We construct the process inductively as follows. Initially, the $N$ particles have configuration $X^N(0)$. Then on $[0,\tau_1)$, drive the particle at location $X_i^N(0)$ by Brownian motion $(W_i(t))_{0\leq s<\tau_1}$, so the locations of the $N$ particles at time $t\in [0,\tau_1)$ are $X_i^N(t):=X_i^N(0)+W_i(t)$ for $i\in [N]$, where $[N]$ denotes the set $\{1,2,\ldots,N\}$. At each stopping time $\tau_i$, we will relabel the particles so that they are ordered. That is $X_1^N(\tau_i)\leq X_2^N(\tau_i)\leq \cdots \leq X_N^N(\tau_i)$ for all $i$, but this ordering may not hold for $t\in (\tau_i,\tau_{i+1})$. We define $\Theta^N:\B{R}^N\to\B{R}^N$ to be the function which orders $\B{R}^N$-valued vectors, and say that $\Theta^N_i(\vec{v})$ is the $i$\textsuperscript{th} smallest element of $\vec{v}$.

Now define $b^N(t):=\inf\left\{x:\Lambda \langle \mu_t^N,\omega(\cdot-x)\rangle \geq f(t)\right\}$, where the notation $\langle \mu,g\rangle$ means $\int_\B{R} g(x)\mu(dx)$. Since $\omega$ is decreasing, $\langle \mu_t^N,1\rangle=1$, and $\Lambda > \|f\|_\infty$% , thus as $x\to\infty$, $\langle \mu_t^N,\omega(\cdot-x)\rangle \to 1$, and as $x\to -\infty$, $\langle \mu_t^N,\omega(\cdot-x)\rangle \to 0$,
, therefore $b^N(t)$ is well defined. Moreover we can observe that if $\omega$ is continuous, then $\frac{\Lambda}{N} \sum_{i=1}^N \omega(X_i^N(t)-b^N(t))$ is exactly $f(t)$. $b^N(t)$ defines a moving boundary. %If $\omega$ is discontinuous, then $\frac{\Lambda}{N}\sum_{i=1}^N \omega(X_i^N(t)-b^N(t))$ differs from $f(t)$ by at most $1/N$ almost surely. 
Then at time $\tau_1$, the particle of rank $I_1$ branches, and simultaneously particle of rank $J_1$ is deleted, where $\B{P}(I_1=i)=1/N$ and independently $\B{P}(J_1=j)=\frac{\omega(\Theta^N_j(X^N(\tau_1-))-b^N(\tau_1-))}{\sum_{\ell=1}^N \omega(\Theta^N_\ell(X^N(\tau_1-))-b^N(\tau_1-))}$%=\frac{\Lambda\omega(\Theta^N_j(X^N(\tau_1-))-b^N(\tau_1))}{Nf(t)}$ 
for $i,j\in [N]$. Repeat this inductively for each interval $[\tau_k,\tau_{k+1}]$. 

%Essentially this process deletes particles, at time $\tau_i$, according to how far to the left of the location $b^N(\tau_i)$. Viewed another way, the particle of rank $j$ is deleted at rate $\Lambda\omega(\Theta^N_j(X^N(t))-b^N(t))$ where $b^N:[0,\infty)\to \B{R}$ is a reference point chosen so that the total branching rate at time $t$ is $Nf(t)$.

\section{Main results}

The main result of this paper is the following hydrodynamic limit, which describes the limiting behaviour of the above described $(\omega,f,\Lambda,N)$-BBM process by the free boundary problem:
\begin{align} \label{omegaNBMMHydro}
    \begin{cases}
        u_t = \frac12 u_{xx} + f(t)u - \Lambda\omega(x-b(t))u & x\in \B{R},\;t>0,\\
        \int_\B{R} u(y,t)dy = 1 & t>0,\\
        u(x,0)=\rho(x) & x\in \B{R}.
    \end{cases}
    \end{align}

\begin{theorem} \label{1stMainThemOmega}
    Let $\omega:\B{R}\to[0,1]$ be a selection weighting such that either $\omega$ is continuous or $\omega(x)=\is_{x\leq 0}$. Let $(\mu_t^N)_{t\geq 0}$ be the measure-valued process describing the $(\omega,f,\Lambda,N)$-BBM process with initial condition $\rho^N$ and let $Q_T^N$ denote its law. Then for any $T\in (0,\infty)$, the sequence of measures $Q_T^N$ has a weak limit, $Q_T^\infty$. Moreover, if $(\mu_t)_{t\in [0,T]}\sim Q_T^\infty$ and $b(t):=\inf\{x:\Lambda\langle \mu_t,\omega(\cdot-x)\rangle \geq f(t)\}$, then $\mu_t$ has density $u(x,t)$ and $(u,b)$ is $Q_T^\infty$-a.s. the unique weak solution to the FBP \eqref{omegaHydroLim}.
\end{theorem}

Using the results of \cite{ettingerHeningEvans}, \cite{ettingerHeningKwong}, we also connect this to the SKIFPP. 


\begin{theorem} \label{2ndMainThemOmega}
    Suppose that $f$ is continuous, and that $\rho\in C^2(\B{R})$ with $\rho$, $\rho'$, and $\rho''$ bounded. Suppose that either $\omega(x)=\is_{x\leq 0}$ or $\omega\in C^3(\B{R})$ is non-increasing and $\{x:\omega(x)\in (0,1)\}$ is bounded. If $(u,b)$ is the unique weak solution to \eqref{omegaNBMMHydro} and $\tau$ is the stopping time:
    \begin{align*}
        \tau:=\inf\left\{t:\Lambda\int_0^t \omega(W_s-b(s))ds\geq U\right\},
    \end{align*}
    where $U\sim Exp(1)$ and $W_0\sim \rho$, then $b$ is the unique barrier such that $\B{P}(\tau>y)=\exp(-\int_0^y f(s)ds)$. That is, $b$ solves the SKIFPP. 
\end{theorem}

\begin{proof}
    %Consider first the case that $\omega(x)=\is_{x\leq 0}$. Then, defining $G(t)=\exp(-\int_0^{t/\Lambda}f(s)ds)$, by Theorem 1.8 \cite{ettingerHeningKwong}, there is a unique weak solution $(v,c)$ solving 
    Let $\tilde{\omega}(x)=\omega(\Lambda^{-1/2}x)$, and suppose that $(v,c)$ is a weak solution to the PDE:
    \begin{align}\label{pdeFromLit}
    \begin{cases}
        v_t(x,t) = \frac12 v_{xx}(x,t) -\tilde{\omega}(x-c(t))v(x,t) & x\in \B{R},\;t>0,\\
        \int_\B{R} v(y,t)dy = G(t):=\exp(-\int_0^{t/\Lambda}f(s)ds) & t>0,\\
        \int_{\B{R}} v(y,0)\omega(y-c(0))dy=f(0)\\
        v(x,0)=\Lambda^{1/2}\rho(\Lambda^{-1/2}x) & x\in \B{R}.\\
    \end{cases}
    \end{align}
    Then it is easily checked that \begin{align*}(\tilde{u}(x,t),\tilde{b}(t)):=\left(\Lambda^{1/2}\exp(\int_0^t f(s)ds)v(\Lambda^{1/2}x,\Lambda t),\Lambda^{-1/2}c(\Lambda t)\right)\end{align*}
    is a weak solution of \eqref{omegaNBMMHydro}, and therefore by assumption coincides with limit $(u,b)$ found in Theorem \ref{1stMainThemOmega}. In the case that $\omega(x)=\is_{x\leq 0}$, Theorem 1.8 \cite{ettingerHeningKwong} tells us that \eqref{pdeFromLit} has a unique weak solution and in the case that $\omega\in C^3(\B{R})$ and $\{x:\omega(x)\in (0,1)\}$ is bounded, Corollary 2.7 \cite{ettingerHeningEvans} tells us that \eqref{pdeFromLit} has a unique (classical) solution. Theorem 1.8 \cite{ettingerHeningKwong} and Corollary 2.7 \cite{ettingerHeningEvans} also tell that $c$ is the unique continuous barrier such that
    \begin{align*}
    \B{P}\left(\inf\left\{t:\int_0^t \omega(W_s-c(s))ds \geq U\right\}>y\right)=G(y)=\exp(-\int_0^{y/\Lambda}f(s)ds).
    \end{align*}
    Then using the fact that $\Lambda^{-1/2}W_{\Lambda z}\overset{d}{=}W_z$ and $c(s)=\Lambda^{1/2}b(s/\Lambda)$, this implies that
    \begin{align*}
        \B{P}\left(\inf\left\{t:\Lambda \int_0^t \omega(W_s - b(s))ds\geq U\right\}>y\right)\exp(-\int_0^y f(s)ds),
    \end{align*}
    which is to say that $b$ solves the SKIFPP. 
\end{proof}


\begin{thebibliography}{99}
\bibitem[]{ettingerHeningEvans} Ettinger, Boris; Evans, Steven N.; Hening, Alexandru. Killed Brownian motion with a prescribed lifetime distribution and models of default. The Annals of Applied Probability 24 (2014) 1-33
\bibitem[]{ettingerHeningKwong} Ettinger, Boris; Hening, Alexandru; Wong, Tak Kwong. The inverse first passage problem for killed Brownian motion. The Annals of Applied Probability 30 (2020) 1251-1275
\end{thebibliography}
\end{document}
